\section{大型稀疏系统与 Krylov 迭代}
\label{sec:08-Numerical-Analysis/05-Numerical-Linear-Algebra/04}

上一节的 Cholesky 把正定结构榨到了尽头：算术量减半，主元不用选，稳定性还是免费的。但它和前面几节共享同一个前提——因子得写得下。

把规模推上去，这个前提就断了。三维网格上的一个稳态扩散问题，离散完有两千万个未知量，矩阵每行七个非零元，压缩存储不到 2\,GB。做一次稀疏 Cholesky 呢？消元会在原本为零的位置填进新的非零元，三维网格上的填充量是 \(O(n^{4/3})\) 量级，因子比原矩阵稠密几个数量级，内存装不下。

还有一类问题连矩阵都没有。训练时要沿自然梯度方向走一步，需要解 \(Fd=g\)，\(F\) 是 Fisher 信息矩阵，维数等于参数量。它从来没有被形成过，也不打算形成：自动微分能以一次前向加一次反向的代价算出 \(Fv\)，仅此而已。你手上不是一个矩阵，是一个函数 \(v\mapsto Fv\)。Hessian 与 Gauss--Newton 矩阵的处境相同。

分解这条路在入口处就关上了。剩下的问题是：只允许调用 \(v\mapsto Av\)，加上向量的加法和内积，能把解逼到多近。

\subsection{反复乘 A 会张出一个越来越好的搜索空间}

从初值 \(x_0\) 和残差 \(r_0=b-Ax_0\) 出发。手头唯一的动作是乘 \(A\)，那就乘：\(r_0,\;Ar_0,\;A^2r_0,\ldots\) 这些向量张成的子空间

\[
\mathcal K_k(A,r_0)=\operatorname{span}\set{r_0,Ar_0,\ldots,A^{k-1}r_0}
\]

叫做 Krylov 子空间。Krylov 方法把近似解限制在仿射空间 \(x_0+\mathcal K_k\) 内，并按某个准则在里面挑最好的那个。

这个约束听上去很随意，写成多项式就不随意了。\(x_k-x_0\) 是 \(r_0\) 的一个次数不超过 \(k-1\) 的矩阵多项式，代进误差 \(e_k=x_\star-x_k\) 得到

\[
e_k=p_k(A)\,e_0,\qquad p_k(0)=1,\ \deg p_k\le k,
\]

其中 \(p_k\) 由算法选定的准则决定。若 \(A\) 可对角化，把 \(e_0\) 按特征方向展开，第 \(i\) 个分量被乘上 \(p_k(\lambda_i)\)。于是收敛问题变成一道逼近问题：\emph{在所有满足 \(p(0)=1\) 的 \(k\) 次多项式里，能不能找到一个在 \(A\) 的全部特征值上都很小的}。

\subsection{共轭梯度在这个子空间里取的是 A-范数最优}

\(A\) 对称正定时，共轭梯度法（CG）在 \(x_0+\mathcal K_k\) 中最小化 \(A\)-范数误差 \(\norm{e}_A=\sqrt{e^{\trans}Ae}\)，而每步只需一次矩阵--向量乘、两次内积和几次向量更新，存三四个向量就够。它不显式构造那组基，也不显式解那个最小化问题——三项递推自动做到了，这是 CG 相对于朴素做法的全部技巧所在。

把上面的多项式论证配上 Chebyshev 多项式，得到经典界

\[
\frac{\norm{e_k}_A}{\norm{e_0}_A}\le 2\left(\frac{\sqrt\kappa-1}{\sqrt\kappa+1}\right)^{k},
\qquad \kappa=\kappa_2(A).
\]

这个界的条件对账值得逐条走一遍：它要求 \(A\) 对称正定；它只用到了区间 \([\lambda_{\min},\lambda_{\max}]\) 这一条信息，把谱在区间内怎么排布全部丢掉了；它是上界，实际收敛只会更好；它推导自精确算术。四条里最要紧的是第二条——被丢掉的那条信息，往往才是决定性的。

注意 \(\sqrt\kappa\) 而不是 \(\kappa\)。条件数一万的系统，梯度下降需要 \(O(10^4)\) 步，CG 的界只要 \(O(10^2)\) 步。这一个平方根，是 CG 值得单独存在的理由。

\subsection{收敛速度写在谱上，不只写在条件数上}

多项式的视角立刻给出一件条件数看不见的事。假设 \(A\) 的特征值聚成三簇，每簇内部很紧。取一个三次多项式，让它的三个根分别落在三簇中心，那么它在所有特征值上都接近零，同时 \(p(0)=1\) 仍可满足。CG 因此在三步左右就把残差压到接近机器精度——不管这三簇之间隔得多远，也就是不管 \(\kappa\) 有多大。

反过来，特征值均匀铺满 \([\lambda_{\min},\lambda_{\max}]\) 时，低次多项式无处安放它的根，只能整体压低，收敛就退回到那个 \(\sqrt\kappa\) 的速率。

\begin{center}
\begin{tikzpicture}[x=0.52cm,y=0.34cm,font=\small,>=stealth]
  \draw[->,black!55] (0,-13.2) -- (0,1.2);
  \draw[->,black!55] (-0.3,-12.5) -- (13.4,-12.5);
  \node[below,font=\footnotesize,black!70] at (6.6,-13.6) {迭代次数 \(k\)};
  \node[rotate=90,font=\footnotesize,black!70] at (-2.6,-6.2) {相对残差};
  \foreach \y/\lab in {0/{\(10^{0}\)},-4/{\(10^{-4}\)},-8/{\(10^{-8}\)},-12/{\(10^{-12}\)}}{
    \draw[black!45] (-0.25,\y) -- (0.25,\y);
    \node[left,font=\scriptsize,black!65] at (-0.3,\y) {\lab};
  }
  \draw[very thick,blue!65] plot[smooth] coordinates
    {(0,0) (1,-1.0) (2,-2.2) (3,-4.0) (4,-6.4) (5,-9.0) (6,-11.6)};
  \node[right,font=\footnotesize,blue!70] at (6.2,-11.7) {谱聚成几簇};
  \draw[very thick,orange!85!black] plot[smooth] coordinates
    {(0,0) (1,-0.35) (2,-0.7) (3,-1.05) (4,-1.5) (5,-1.9) (6,-2.35)
     (7,-2.8) (8,-3.3) (9,-3.8) (10,-4.3) (11,-4.85) (12,-5.4)};
  \node[right,font=\footnotesize,orange!85!black] at (12.1,-5.5) {谱铺开};
  \begin{scope}[shift={(0,-17.0)}]
    \draw[black!55] (0,0) -- (11.4,0);
    \foreach \x in {0.8,1.05,1.3,5.3,5.55,9.7,9.95,10.2,10.45}{
      \fill[blue!65] (\x,0) circle (1.7pt);}
    \node[right,font=\footnotesize,black!70] at (11.6,0) {聚集的谱};
  \end{scope}
  \begin{scope}[shift={(0,-20.4)}]
    \draw[black!55] (0,0) -- (11.4,0);
    \foreach \x in {0.7,1.5,2.3,3.1,3.9,4.7,5.5,6.3,7.1,7.9,8.7,9.5,10.3}{
      \fill[orange!85!black] (\x,0) circle (1.7pt);}
    \node[right,font=\footnotesize,black!70] at (11.6,0) {铺开的谱};
  \end{scope}
\end{tikzpicture}

\emph{两个系统的 \(\lambda_{\min}\) 与 \(\lambda_{\max}\) 一样，条件数因此也一样。谱聚成几簇时，一个低次多项式就能同时压住所有簇，残差成阶梯状掉下去；谱铺满整个区间时，同样的迭代次数只换来两三个数量级。}
\end{center}

图里那条蓝线不是特例。离散化良好的椭圆型算子、加了合适预条件的鞍点系统、低秩扰动过的单位阵，谱都长成簇状。判断一个迭代求解器会不会成功，先问谱长什么样，而不是先问 \(n\) 有多大。

\subsection{预条件改造的是谱，不是解}

既然速度由谱决定，加速的办法就明确了：换一个谱更好的等价系统。找一个容易求逆的 \(M\approx A\)，改解 \(M^{-1}Ax=M^{-1}b\)。解一个字都没变，变的是迭代看到的谱——理想情况下 \(M^{-1}A\) 的特征值挤在 1 附近。

好的预条件器要同时满足几件互相拉扯的事。它每次应用必须便宜，否则每步成本超过省下的步数；它得把谱真的聚拢，而不只是把条件数从 \(10^{8}\) 降到 \(10^{7}\)；用 CG 时它必须保持对称正定，否则推导所依赖的内积不再是内积；在并行机器上它还不能引入长串的串行依赖。常见的选择沿着``便宜''到``有效''这条轴排开：对角缩放几乎不花钱、效果也有限；不完全 LU 或不完全 Cholesky 丢掉填充量大的位置，是通用的中间档；多重网格对椭圆型问题能把迭代次数压到与 \(n\) 无关，要付出的是针对具体问题做一遍设计。

有一个反直觉但常见的结论：迭代次数最少的预条件器往往不是最快的。比较的对象应该是总时间，不是迭代计数。

\subsection{CG、MINRES 与 GMRES 不能互相替换}

CG 要求对称正定。把一个不定矩阵交给 CG，\(\norm{\cdot}_A\) 不再是范数，三项递推赖以成立的正交性也塌了，结果不是``慢一点''，而是可能停滞或发散，而残差曲线看上去还挺正常。

对称但不定的矩阵用 MINRES：它保留短递推的低存储优势，改成最小化二范数残差。一般非对称矩阵用 GMRES，它在 Krylov 子空间中最小化残差，代价是必须保存并正交化全部基向量，第 \(k\) 步的存储和正交化成本都随 \(k\) 线性增长。实践中用重启版本 GMRES(\(m\))，把内存钉死在 \(m\) 个向量，但重启会丢掉已积累的子空间信息，可能显著变慢甚至停滞。BiCGSTAB 一类方法用短递推换掉最优性，通常更省内存，收敛曲线却不再单调。

选择因此有一条清晰的主线：先看矩阵的对称性和定性，再看内存预算，最后才比较收敛曲线。

\subsection{残差小仍然不等于解对}

迭代求解器的停止准则通常是相对残差 \(\norm{r_k}/\norm{b}\) 降到某个容差以下。这个量便宜、可算、单调，但它衡量的是后向误差，不是前向误差——\hyperref[sec:08-Numerical-Analysis/01-Floating-Point-and-Error/04]{``前向误差、后向误差与稳定算法''}里那笔账在这里原样重演：\(x_\star-x_k=A^{-1}r_k\)，从残差推到误差要乘一次 \(\kappa(A)\)。病态系统上停在 \(10^{-10}\) 的残差，解可能只有几位有效数字。

还有两个容易踩的坑。预条件后监控的往往是 \(M^{-1}r_k\) 而非 \(r_k\)，两者的量级可以差很多，容差要按哪一个设需要写清楚。另一个是浮点下 Krylov 基会逐渐失去正交性，教科书里``精确算术下至多 \(n\) 步收敛''这句话在实际计算中不成立，不能当作兜底保证。求解器停滞时，先分别排查预条件失效、尺度不当、右端落在近似零空间方向上这几种原因，再考虑换算法。

\subsection{直接法看规模，迭代法看谱}

这一节和前三节的分界可以收成一句话：直接法的可行性由规模和稀疏结构决定，一旦因子存得下、算得完，精度就有后向稳定性兜底；迭代法的可行性由谱决定，存储从来不是它的瓶颈，收敛才是。规模不大时不要为了``现代''而上迭代法，因为直接法给的保证更硬；规模上去之后，能不能算完的问题就转化成了谱好不好的问题，而预条件是唯一能动这个变量的旋钮。

Krylov 子空间还能做另一件事。\(\set{r_0,Ar_0,A^2r_0,\ldots}\) 里积累的不只是解的信息，还有谱的信息——Lanczos 与 Arnoldi 正是从这组基里提取特征值。下一节从最朴素的做法开始：只反复乘矩阵，看会浮出什么。
